\documentclass[11pt,a4paper]{article}

% --- Pacotes de Idioma e Codificação ---
\usepackage[utf8]{inputenc}
\usepackage[T1]{fontenc}
\usepackage[portuguese]{babel}

% --- Design e Layout ---
\usepackage[margin=2.5cm]{geometry}
\usepackage{charter} % Fonte profissional e legível
\usepackage{lastpage}
\usepackage{fancyhdr}
\usepackage[table]{xcolor}
\usepackage{tabularx}
\usepackage{booktabs}
\usepackage{xcolor}
\usepackage{enumitem}
\usepackage{amssymb}

% --- Configuração de Cores ---
\definecolor{primaryblue}{RGB}{0, 51, 102}

% --- Cabeçalho e Rodapé ---
\pagestyle{fancy}
\fancyhf{}
\fancyhead[L]{\small \textbf{Registo de Entrevista} | EcoRide Solutions}
\fancyhead[R]{\small \thepage\ de \pageref{LastPage}}

\renewcommand{\headrulewidth}{0.4pt}

% --- Ambiente Personalizado para Transcrição ---
% Facilita a distinção entre Entrevistador (E) e Entrevistado (R)
\newlist{transcription}{description}{1}
\setlist[transcription]{
	labelsep=1em,
	leftmargin=3.5em,
	style=nextline,
	font=\bfseries\color{primaryblue}
}

\newcommand{\pergunta}[1]{\item[P:] \textit{#1}}
\newcommand{\resposta}[1]{\item[R:] #1}

% --- Início do Documento ---
\begin{document}
	
	% Título Principal
	{\noindent\LARGE\bfseries\color{primaryblue} Relatório de Entrevista Individual}
	\\[0.5cm]
	
	% --- Tabela de Metadados ---
	\renewcommand{\arraystretch}{1.5}
	\begin{tabularx}{\textwidth}{@{}|l|X|l|X|@{}}
		\hline
		\rowcolor{gray!10} \textbf{ID Entrevista} & \#001 & \textbf{Data} & 26 de Fevereiro de 2026 \\ \hline
		\textbf{Hora Início} & 10:00 & \textbf{Hora Fim} & 10:30 \\ \hline
		\textbf{Local/Meio} & Sala de Reuniões B & \textbf{Estado} & Finalizada / Transcrita \\ \hline
	\end{tabularx}
	
	\vspace{0.5cm}
	
	% --- Participantes ---
	\begin{tabularx}{\textwidth}{@{}|l|X|@{}}
		\hline
		\rowcolor{gray!10} \textbf{Intervenientes} & \textbf{Cargo / Função} \\ \hline
		\textbf{Entrevistador(a)} & Luís Soares \\ \hline
		\textbf{Entrevistado(a)}  & João Martins - Dono e Gestor do Armazém da oficina \\ \hline
	\end{tabularx}
	
	\vspace{0.8cm}
	
	\hrule % Linha separadora opcional
	\vspace{0.2cm}
	
	% --- Transcrição ---
	\section*{1. Transcrição / Notas da Entrevista}
	
	\begin{transcription}
		
		\pergunta{Que responsabilidades detém o seu papel de gerente nesta empresa ?}
		
		\resposta{Como gerente da oficina, o meu papel passa sobretudo por garantir que tudo funciona de forma organizada e sustentável, tanto no dia a dia como a longo prazo. Sou responsável por acompanhar o volume de trabalho, assegurar que os prazos são cumpridos, controlar a qualidade das reparações e garantir que os clientes saem satisfeitos. Também tenho de manter a empresa financeiramente saudável, o que implica acompanhar custos, margens e resultados mensais. Além disso, trato da gestão do armazém e da logística, porque é essencial garantir que temos sempre as peças necessárias e que o fluxo de entrada e saída das trotinetes decorre sem confusões. No fundo, o meu papel é coordenar todas as áreas para que a oficina funcione de forma eficiente e sem interrupções.}
		
		\pergunta{Para além do nome e contacto, que informação é guardada relativamente aos clientes? Além da marca e modelo, que detalhes técnicos são vitais para a descrição de uma trotinete? Podemos ter várias trotinetes associadas ao mesmo cliente?}
		\resposta{Relativamente aos clientes, para além do nome e contacto, costumamos guardar também o número de contribuinte, porque é necessário para a faturação. No que toca às trotinetes, além da marca e modelo, é essencial termos o número de
		série, porque é isso que nos permite identificar exatamente cada unidade. Também é
		importante sabermos o estado geral da trotinete à entrada, o tipo de motor, o tipo de bateria, e se tem algum acessório específico, como por exemplo, uma luz ou assento. Sim, é comum um cliente ter mais do que uma trotinete por exemplo, empresas de micromobilidade ou clientes que têm várias em casa.}

		\pergunta{Quando contrata alguém para trabalhar na empresa, que dados pede-lhe ? Como é que calcula o valor do salário para lhe pagar ao fim do mês ?}
		\resposta{Sempre que contrato alguém, preciso de saber alguns detalhes pessoais, como nome, morada, número de telemóvel, email, data de nascimento, o NISS, o NIF, o NUS e o seu IBAN. Quando estabelecemos o contrato, é acordado um valor e ao fim do mês faço a transferência para o IBAN associado ao funcionário.
		}
		
		\pergunta{Existe a possibilidade de alterar o valor do salário pago ao funcionário ? Costuma pagar por horas extraordinárias ?}
		\resposta{Com o decorrer da carreira de alguns dos funcionários que já trabalharam na oficina, já tive de realizar alguns aumentos e por isso, sim, posso precisar de alterar o valor pago. Em relação às horas extraordinárias, existem alturas em que o fluxo de trabalho aumenta e por isso é necessário que o pessoal faça horas extra, e, por isso, esse valor é acrescentado ao salário acordado.}
		
		\pergunta{Acha que seria relevante ter o valor de salário a pagar ao funcionário ao fim do mês caso ele tivesse feito horas extraordinárias?}
		\resposta{Sim, quando eles fazem horas, eu preciso de manter um registo para que ao fim do mês eu tenha uma noção do valor extra a transferir. Apesar de eu manter este registo, já houveram situações em que enganei-me a transferir este valor e tive de realizar correções.}
		
		\pergunta{Quando um cliente entra na loja com uma trotinete avariada, o que é que é anotado primeiro?}
		\resposta{Quando um cliente entra na loja com uma trotinete avariada, o primeiro passo é recolher os dados básicos do cliente e da própria trotinete. Normalmente, a secretária anota o nome do cliente, o email e o número de telemóvel, e faz uma descrição rápida do problema que o cliente relata. Também se regista a marca e o modelo da trotinete, e se possível, o número de série. Este processo é feito em papel, geralmente numa ficha ou folha própria para isso, que depois é colocada numa pasta com outras ordens.}
		
		\pergunta{Em relação à abertura de uma ordem de serviço, em que estado a trotinete fica? Como é que o mecânico sabe que tem uma trotinete nova à espera?}
		\resposta{Assim que a trotinete entra, consideramos que está por diagnosticar ou pendente de avaliação. A secretária costuma avisar verbalmente os mecânicos ou deixar a ficha numa	zona onde eles sabem que há novas entradas. Não há um sistema automático para isso os mecânicos vão vendo ao longo do dia se há novas trotinetes na zona de entrada ou se há fichas novas na pasta.}
		
		\pergunta{No diagnóstico, é costume dar um orçamento fixo ou o valor pode mudar a meio da reparação? Como é que o cliente é informado dessa mudança?}
		\resposta{No diagnóstico, tentamos sempre dar um orçamento ao cliente antes de avançar. No entanto, nem sempre conseguimos prever tudo à primeira vista. Às vezes, ao desmontar a trotinete, descobrimos outros problemas que não estavam visíveis no início. Nesses casos, a secretária entra em contacto com o cliente, normalmente por telefone ou email, para explicar a situação e confirmar se ele quer avançar com os custos adicionais. Só depois de termos essa confirmação é que continuamos.}
		
		\pergunta{Relativamente ao processo de reparação, o mecânico pode começar a reparar sem o cliente aprovar o orçamento? Como é que o sistema deve travar isso?}
		\resposta{O mecânico não deve começar a reparação sem o cliente aprovar o orçamento. Isso é uma regra da casa, até para nos protegermos de mal-entendidos. Mas como não temos um sistema que bloqueie isso, dependemos da comunicação entre a receção e os mecânicos. Pode acontecer, em dias mais atarefados, que alguém comece a mexer numa trotinete antes da aprovação estar formalizada, o que não é o ideal. Por isso, seria importante termos uma forma mais segura de garantir que só se avança quando o cliente já aprovou.
		}
		
		\pergunta{Para gerir a empresa e poder acompanhar o estado das ordens de serviço pendentes, o que é que faz ?}
		\resposta{Para acompanhar o estado das reparações, dependo muito das fichas em papel e da comunicação com a secretária e com os mecânicos. Vou vendo as ordens que estão na bancada, pergunto como estão a avançar e confirmo se já houve resposta dos clientes quando há orçamentos pendentes. Não tenho um sistema centralizado, por isso acabo por andar fisicamente pela oficina a ver o que está parado, o que está em reparação e o que já está pronto. É um processo muito manual, mas é a única forma de garantir que nada fica esquecido.}

		\pergunta{Como é que sabe os lucros e despesas ao longo do mês ?}
		\resposta{No final do mês, junto as faturas emitidas, os valores recebidos e os custos principais, como compras de peças, salários e despesas fixas. Faço isso numa folha de Excel onde lanço tudo manualmente. Não é muito sofisticado, mas dá-me uma ideia geral de como correu o mês. O que me interessa é perceber se estamos a crescer, se estamos a gastar mais do que devíamos e se as reparações estão a dar a margem que espero.}

		\pergunta{As cobranças dos serviços prestados são decididas por você ? Existe a possibilidade de querer alterar esses valores ?}
		\resposta{Sim, sou eu que decido os preços dos serviços e das peças. A secretária e os mecânicos não têm autonomia para alterar valores, porque sou eu que acompanho os preços de compra, as margens e o que é praticado no mercado. Às vezes preciso de ajustar um preço, seja porque o custo da peça mudou, seja porque a reparação foi mais complexa do que o habitual. Por isso, sim, é importante que eu tenha a possibilidade de alterar valores quando for necessário, mas sempre de forma controlada.}
		
		\pergunta{Para além da gestão da empresa, o que é que faz para controlar as peças de reparo que são compradas e gastas ?}
		\resposta{Para controlar as peças de reparação, acabo por ter de acompanhar tudo muito de perto, porque o armazém é uma parte crítica do negócio. Sempre que chega material novo, sou eu que confiro as encomendas, verifico se as quantidades estão corretas e se os preços batem certo com o que foi combinado com o fornecedor. Depois disso, registo a entrada
		num caderno ou numa folha de cálculo, onde tenho uma lista com as peças, as
		quantidades e, o preço de compra. Finalmente, organizo as peças no armazém para que os mecânicos as encontrem facilmente. 
		
		Ao longo do dia, vou controlando o que sai, normalmente através das fichas das ordens de serviço, onde anoto as peças que os mecânicos dizem ter usado. No final da semana, faço uma revisão geral para perceber o que foi gasto, o que está a acabar e o que precisa de ser encomendado novamente. Tento manter sempre um equilíbrio entre ter stock suficiente para não parar as reparações e não acumular peças que demoram a sair, porque isso representa dinheiro parado.}
		
		\pergunta{Como é que são identificadas as peças ? Usa alguma referência ?}
		\resposta{Para identificar as peças, uso sempre as referências dos fornecedores, como são sempre únicas, caso volte a precisar de encomendar, posso usar a referência do fornecedor para facilitar-lhe a identificação da peça que preciso. Para além disto, no caso de peças mais caras, como baterias ou motores, tento sempre apontar o número de série. Isto é importante para sabermos exatamente qual foi a unidade
		instalada em cada trotinete, principalmente se houver uma reclamação ou se for preciso acionar a garantia.}
		
		\pergunta{Anteriormente mencionou o facto de que quando os mecânicos precisam de uma peça, apontam na folha da ordem de serviço em que a peça foi usada. Não existe a possibilidade de em dias mais atarefados eles esquecerem-se disso ?}
		\resposta{Quando uma peça é usada numa reparação, o normal é o mecânico vir ter comigo e pedir
		a peça. Eu entrego e registo logo na folha da ordem de serviço, para garantir que não
		me esqueço. Mas há dias mais complicados em que o mecânico vai buscar diretamente
		e só me avisa depois. Nesses casos, tento confirmar o mais depressa possível e atualizar
		o registo, mas admito que às vezes pode escapar alguma coisa.}
		
		\pergunta{E se um mecânico levantar uma peça para um diagnóstico e depois não a usar? Como é que
		faz o processo de devolução da peça ao stock?}
		\resposta{Se o mecânico levantar uma peça para testar durante o diagnóstico e depois não a usar,
		costumo pedir que me devolva logo. Quando isso acontece, volto a colocá-la no sítio e
		atualizo o stock, como se fosse uma entrada. Para além disso, removo-a da lista de pelas usadas da ficha da ordem de serviço respetiva.}
		
		\pergunta{Como é que sabe que está na hora de encomendar mais material?}
		\resposta{Quando preciso de saber se está na hora de encomendar mais material, costumo verificar
		o stock manualmente ou através das anotações que mantenho atualizadas. Tenho uma
		ideia dos mínimos que gosto de ter para cada tipo de peça, mas nem sempre é fácil
		controlar tudo ao detalhe. }
	
		\pergunta{Guarda os contactos dos fornecedores ?  Seria útil conseguir gerar uma lista de
		compras automática para enviar por e-mail ao fornecedor ?}
		
		\resposta{Sim, guardo os contactos dos fornecedores num ficheiro Excel, com os nomes, telefones e
		e-mails. Também tenho lá os preços habituais e os prazos de entrega. Geralmente, quando preciso de fazer uma encomenda, ligo ao fornecedor e digo-lhe a referência para saber se ele tem a pela necessária. Outras vezes, mando uma lista por email com as peças que preciso ou escrevo uma lista para ir a uma loja que tenha perto.}
	
		\pergunta{Como controla as peças que enviou para trás por estarem com defeito ? }
		\resposta{Quando tenho de devolver peças com defeito, costumo separá-las numa caixa marcada
		como “para devolução” e aponto num papel o que está lá dentro, com a data e o motivo.
		Às vezes perco o papel e esqueço-me de algumas das pelas devolvidas, principalmente se o fornecedor demorar a responder, o que gera alguma confusão quando sou contactado por ele.}
		
	\end{transcription}
	\vspace{0.5cm}

	\hrule % Linha separadora opcional
	\vspace{0.2cm}
	
	\section*{2. Requisitos obtidos}
	
		\begin{tabularx}{\textwidth}{@{}|l|X|@{}}
			\hline
			\rowcolor{gray!10} \textbf{Número} & \textbf{Descrição} \\ \hline
			\textbf{1} & O sistema deve disponibilizar ao gerente uma interface centralizada que permita visualizar, em tempo real, o estado de todas as ordens de serviço (pendentes, em reparação, concluídas e aguardando aprovação). Deve permitir a filtragem por estado, data, cliente ou mecânico responsável. \\ \hline
			\textbf{2}  & O sistema deve permitir registar digitalmente todas as entradas de peças (com quantidade, fornecedor, preço, data e, opcionalmente, número de série). O registo deve ser obrigatório no momento da movimentação para garantir rastreabilidade. \\ \hline
			\textbf{3}  & O sistema deve permitir associar peças utilizadas diretamente à ordem de serviço correspondente, garantindo que cada peça retirada do stock é registada e descontada automaticamente. \\ \hline
			\textbf{4}  & O sistema deve permitir registar devoluções de peças não utilizadas, restaurando automaticamente o stock e removendo a peça da ordem de serviço onde tinha sido inicialmente registada. \\ \hline
			\textbf{5}  & O sistema deve permitir configurar níveis mínimos por peça e gerar alertas automáticos quando o stock atingir ou ultrapassar esses limites. \\ \hline
			\textbf{6}  & O sistema deve gerar automaticamente listas de peças a encomendar com base nos níveis mínimos, consumo recente e stock atual, permitindo exportação para uma folha ou envio direto por e-mail ao fornecedor. \\ \hline
			\textbf{7}  & O sistema deve permitir armazenar e consultar contactos de fornecedores, incluindo nome, telefone, e-mail, preços habituais e prazos de entrega. \\ \hline
			\textbf{8}  & O sistema deve permitir registar peças devolvidas ao fornecedor, incluindo motivo, data, estado da devolução e histórico de comunicação. \\ \hline
			\textbf{9}  & O sistema deve permitir ao gerente atualizar preços de serviços e peças, mantendo histórico de alterações e garantindo que apenas perfis autorizados podem modificar valores. \\ \hline
			\textbf{10}  & O sistema deve gerar relatórios mensais automáticos com receitas, despesas, margens e indicadores financeiros. \\ \hline
		\end{tabularx}


\end{document}